\subsection{Multiplication by the coordinate and the prime form}
\label{subsec:fd-prime-commutator}

There is an exact arithmetic form identity for physical
multiplication by the coordinate. It also identifies a domain issue
in an attempted comparison of local multiplicity traces with the
ordinary spectral trace.

Let \(D=-i\partial_x\) and \(Xf(x)=xf(x)\), initially on compact
smooth functions. For \(f,g\) in this class write
\(h=f*\widetilde g\), with
\(\widetilde g(x)=\overline{g(-x)}\), and let \(W\) be the
arithmetic Weil functional. Direct convolution gives
\[
 (Xf)*\widetilde g-f*\widetilde{Xg}=xh(x).
\]
Consequently
\begin{align}
 q(Xf,g)-q(f,Xg)
 &=W(xh)\notag\\
 &=\int_{\mathbb R}xh(x)(e^{x/2}+e^{-x/2})\,dx\notag\\
 &\quad-\int_0^\infty
 \frac{xe^{-x/2}(h(x)-h(-x))}{1-e^{-2x}}\,dx\notag\\
 &\quad-\sum_{n\ge2}\frac{\Lambda(n)\log n}{\sqrt n}
                  \bigl(h(\log n)-h(-\log n)\bigr).
 \label{eq:fd-prime-commutator}
\end{align}
The prime sum is finite because \(h\) has compact support.
The local term in the Weil formula vanishes since \(xh(x)\) is zero
at the origin. The remaining integral is regular there because
\(h(x)-h(-x)=O(x)\). On the spectral side the same identity is
\[
 q(Xf,g)-q(f,Xg)
 =i\sum_\gamma m_\gamma
        (\widehat f\,\overline{\widehat g})'(\gamma),
\]
as \(\widehat{Xf}=i\widehat f'\). Also
\([D,X]f=-if\) holds exactly on this test space.

\begin{proposition}[A nonzero residual for global multiplication]
\label{prop:fd-multiplication-domain}
The operator \(X\), defined on compact smooth mean-zero tests and
viewed as an operator into the full Weil completion, is not closable.
For every signed ordinate \(\gamma\) there are such tests \(f_R\)
with
\[
 \|f_R\|_q\longrightarrow0,
 \qquad
 Xf_R\longrightarrow i\mathbf1_{\{\gamma\}}
       \quad\hbox{in }\ell^2(\Gamma,m).
\]
The squared norm of the limiting vector is \(m_\gamma\).
\end{proposition}

\begin{proof}
Choose an even \(\varphi\in C_c^\infty(\mathbb R)\) of integral
one, and put \(K=\widehat\varphi\). Before correcting the mean,
define
\[
 g_R(x)=R^{-1}e^{i\gamma x}\varphi(x/R),\qquad
 f_R^{\,0}=(D-\gamma)g_R.
\]
Their Fourier transforms satisfy
\[
 F_R^{\,0}(z)=(z-\gamma)K(R(z-\gamma)),
\]
\[
 \widehat{Xf_R^{\,0}}(z)
 =i\bigl(K(R(z-\gamma))
        +R(z-\gamma)K'(R(z-\gamma))\bigr).
\]
At \(z=\gamma\) these values are zero and \(i\), respectively.
Every distinct ordinate is isolated, and the counted zero bound gives
\[
 S_\gamma=\sum_{\lambda\ne\gamma}
             \frac{m_\lambda}{(\lambda-\gamma)^2}<\infty.
\]
Indeed the sum is finite locally after removing the isolated
ordinate, and the tail converges by the unit-interval count
\(O(\log(|\lambda|+3))\).

Both \(K\) and \(K+vK'\) are Schwartz functions. Hence, for
\(R\ge1\) and \(d\ne0\),
\[
 |dK(Rd)|\le C R^{-2}|d|^{-1},
 \qquad
 |K(Rd)+RdK'(Rd)|\le C R^{-1}|d|^{-1}.
\]
The first estimate follows from the boundedness of \(v^2K(v)\),
and the second from the boundedness of \(v(K(v)+vK'(v))\).
Summing their squares against the zero weights proves
\[
 q(f_R^{\,0})\le C R^{-4}S_\gamma\longrightarrow0,
\]
\[
 \|Xf_R^{\,0}-i\mathbf1_{\{\gamma\}}\|_q^2
 \le C R^{-2}S_\gamma\longrightarrow0.
\]

For the mean correction choose one fixed
\(b\in C_c^\infty(\mathbb R)\) with \(\int b=1\), and set
\[
 f_R=f_R^{\,0}-F_R^{\,0}(0)b.
\]
There is no zero ordinate, so
\(F_R^{\,0}(0)=-\gamma K(-R\gamma)\to0\).
Both \(b\) and \(Xb\) have finite Weil norm. The correction
therefore vanishes in both displayed limits and makes
\(\int f_R=0\). The limiting graph pair is
\((0,i\mathbf1_{\{\gamma\}})\), whose second entry is nonzero.
This is exactly the obstruction to closability.
The same sequence can have its first physical moment zero as well.
Indeed \((F_R^{\,0})'(0)\to0\). Choose fixed compact smooth
functions \(b_0,b_1\) with moment pairs
\((\int b_j,\int xb_j)=(1,0),(0,1)\), respectively, and subtract
the two moments of \(f_R^{\,0}\) times these functions. Both
coefficients tend to zero, so both graph limits remain unchanged.
Then \(f_R\) and \(Xf_R\) are both mean-zero tests. For example,
one may take an even mean-one \(b_0\) and \(b_1=-b_0'\).
\end{proof}

The trace comparison can be stated without identifying two different
trace functionals. The unitary map
\[
 U:\ell^2(\Gamma,m)\longrightarrow\ell^2(\Gamma),
 \qquad (Ua)_\gamma=\sqrt{m_\gamma}\,a_\gamma
\]
conjugates every diagonal translation
\(a_\gamma\mapsto e^{-it\gamma}a_\gamma\), and every defined
diagonal function of \(D\), to the same diagonal operation. In
particular, for a spectral function \(v\ge0\) with
\(\sum_\gamma m_\gamma v(\gamma)<\infty\),
the ordinary Hilbert-space trace is
\[
 \operatorname{tr}v(D)=\sum_{\gamma\in\Gamma}v(\gamma),
\]
whereas the zero-divisor functional is
\[
 \tau_{\rm div}(v)=\sum_{\gamma\in\Gamma}m_\gamma v(\gamma).
\]
Their exact difference is
\(\sum_\gamma(m_\gamma-1)v(\gamma)\); for
\(v(\gamma)=\gamma^{-2k}\), it is the limiting trace defect
already proved above. Thus conjugating the commuting translation
operators does not change their ordinary trace or manufacture this
additional term.

The form commutator \eqref{eq:fd-prime-commutator} retains the
derivative evaluations and hence additional information about the
physical embedding. The proposition shows why passing directly from
that identity to a cyclic trace of a global closed multiplication
operator is invalid: the stated test operator has the explicit
nonzero graph residual above. A bound controlling these derivative
channels in a suitable growing-window trace comparison has not been
established here. The prime-translation identity remains exact, and
no conclusion about the impossibility of an arithmetic comparison
is drawn from this failed operator passage.
